\section*{Chapitre \ref{ch:g12:light-as-wave} --- La lumière comme onde~: diffraction et interférences}

\begin{solution}{exo:g12:light-as-wave:1}
$f = c/\lambda$~: rouge $\num{3.00e8}/\num{7.0e-7} \approx
\qty{4.3e14}{Hz}$~; violet $\approx \qty{7.5e14}{Hz}$ --- le violet est
plus élevé (\omterm{def:g12:mechanical-waves:wavelength}{longueur d'onde} plus courte, même célérité $c$).
\end{solution}

\begin{solution}{exo:g12:light-as-wave:2}
$\theta = \num{5.32e-7}/\num{2.0e-4} = \qty{2.7e-3}{rad}$~;
$L = 2\theta D = 2 \times \num{2.66e-3} \times 1.5 \approx
\qty{8.0}{mm}$.
\end{solution}

\begin{solution}{exo:g12:light-as-wave:3}
Voix~: $\lambda = 340/425 = \qty{0.80}{m}$, donc $\lambda/a = 1.0$ ---
\omterm{def:g12:light-as-wave:diffraction}{diffraction} maximale, le \omterm{def:g10:signals-and-waves:sound}{son} inonde le couloir.
Lumière~: $\lambda/a \approx \num{5.5e-7}/0.80 \approx \num{7e-7}$ ---
étalement négligeable~: la lumière reste droite, donc le coin la cache.
\end{solution}

\begin{solution}{exo:g12:light-as-wave:4}
En $M$~: $\delta = \qty{6.0}{cm} = 3\lambda$, \omterm{prop:g12:light-as-wave:conditions}{constructive} --- fort.
En $N$~: $\delta = \qty{5.0}{cm} = (2 + \tfrac12)\lambda$, \omterm{prop:g12:light-as-wave:conditions}{destructive}
--- (presque) silencieux.
\end{solution}

\begin{solution}{exo:g12:light-as-wave:5}
$i = \lambda D/a = \num{6.33e-7} \times 1.8/\num{2.5e-4} \approx
\qty{4.6}{mm}$. Dans $\pm\qty{1.0}{cm}$~: $|k| \leq 10/4.6 = 2.2$,
donc $k = -2, \dots, 2$~: cinq franges brillantes.
\end{solution}

\begin{solution}{exo:g12:light-as-wave:6}
$d = 2\lambda D/L = 2 \times \num{6.5e-7} \times 1.6/0.028 \approx
\qty{74}{\micro m}$.
\end{solution}

\begin{solution}{exo:g12:light-as-wave:7}
$i = \qty{4.7}{mm}$, donc $\lambda = i\,a/D = \num{4.7e-3} \times
\num{2.0e-4}/1.5 \approx \qty{6.3e-7}{m} = \qty{630}{nm}$~: rouge.
Mesurer dix \omterm{def:g12:light-as-wave:fringes}{interfranges} divise l'erreur de lecture de la règle
par dix.
\end{solution}

\begin{solution}{exo:g12:light-as-wave:8}
Les lampes ont (largement) des fréquences assorties mais ne sont
\emph{pas en phase}~: chacune émet de courts trains d'\omterm{def:g10:signals-and-waves:wave}{ondes} à phase
sautant aléatoirement, si bien que les sources sont incohérentes et que
les franges instantanées se déplacent des milliards de fois par seconde
--- l'œil voit la moyenne, uniforme. Young éclaire les deux fentes
depuis \emph{une} \omterm{def:g10:signals-and-waves:wave}{onde}~: les deux copies héritent de chaque saut de
phase ensemble et restent en phase.
\end{solution}

\begin{solution}{exo:g12:light-as-wave:9}
De la première à la septième $= 6i = \qty{18.0}{mm}$, donc
$i = \qty{3.0}{mm}$~;
$\lambda = i\,a/D = \num{3.0e-3} \times \num{3.0e-4}/1.4 \approx
\qty{6.4e-7}{m} \approx \qty{640}{nm}$.
\end{solution}

\begin{solution}{exo:g12:light-as-wave:10}
$\sin\theta = k\lambda/d = 0.396\,k$~: $k = 1$, $\theta = 23^\circ$~;
$k = 2$, $\sin\theta = 0.79$, $\theta = 52^\circ$~; $k = 3$ exigerait
$\sin\theta = 1.19 > 1$~: impossible. Ordre maximal $k = 2$.
\end{solution}

\begin{solution}{exo:g12:light-as-wave:11}
$\theta = \num{5.5e-7}/0.10 = \qty{5.5e-6}{rad}$. Les phares
sous-tendent $1.5/\num{1.0e5} = \qty{1.5e-5}{rad}$, environ trois fois
le flou~: séparés, tout juste.
\end{solution}

\begin{solution}{exo:g12:light-as-wave:12}
Les réflexions sur les faces avant et arrière \omterm{prop:g12:light-as-wave:conditions}{interfèrent}~; en
\omterm{def:g10:light-spectra:white}{lumière blanche} chaque épaisseur renforce certaines
\omterm{def:g12:mechanical-waves:wavelength}{longueurs d'onde} et en annule d'autres, d'où une couleur.
L'épaisseur du film ne dépend que de la hauteur, donc les points de
même couleur forment des bandes horizontales. L'écoulement
\omterm{def:g11:lenses-and-eye:lens}{amincit} le film, si bien que l'épaisseur qui peignait une
couleur donnée siège toujours plus bas~: les bandes rampent vers le bas.
\end{solution}

\begin{solution}{exo:g12:light-as-wave:13}
$i = \lambda D/a$~: violet $\num{4.0e-7} \times 1.5/\num{2.0e-4} =
\qty{3.0}{mm}$~; rouge $\qty{5.6}{mm}$. Centre~: toutes les couleurs
brillantes --- une frange blanche. À quelques millimètres les motifs
colorés se sont décalés~: franges irisées. Plus loin, les maxima de
toutes les couleurs se chevauchent partout~: flou blanc uniforme.
\end{solution}

\begin{solution}{exo:g12:light-as-wave:14}
$\theta = \num{5.5e-7}/2.4 = \qty{2.3e-7}{rad}$~; à \qty{250}{km}
cela fait $\num{2.3e-7} \times \num{2.5e5} \approx \qty{6}{cm}$. Il
peut compter des voitures et repérer une personne, mais le journal
(lettres millimétriques) est vingt fois sous la limite de
\omterm{def:g12:light-as-wave:diffraction}{diffraction}~: l'affirmation est un mythe.
\end{solution}

\begin{solution}{exo:g12:light-as-wave:15}
La \omterm{def:g10:signals-and-waves:frequency}{fréquence} est fixée par la source~: inchangée. La célérité
tombe à $c/n$, donc $\lambda' = \lambda/n$ et
$i' = i/n = 4.0/1.33 \approx \qty{3.0}{mm}$~: les franges se resserrent.
\end{solution}

\begin{solution}{pb:g12:light-as-wave:1}
\textbf{1.} $\theta = \lambda/a = \num{6.33e-7}/\num{1.00e-4} =
\qty{6.3e-3}{rad}$.

\textbf{2.} Demi-largeur $D\tan\theta \approx D\theta = \lambda D/a$,
donc $L = 2\lambda D/a = 2 \times \num{6.33e-7} \times
3.00/\num{1.00e-4} = \qty{3.8}{cm}$.

\textbf{3.} $(3.8 - 3.7)/3.8 \approx 3\%$~: dans la précision d'une
règle --- validé.

\textbf{4.} $L \propto 1/a$~: diviser $a$ par deux double $L$, donc
$L = \qty{7.6}{cm}$. Obstacle ou fente plus étroit, motif plus large.

\textbf{5.} $\tan(\num{6.33e-3}) = \num{6.33008e-3}$~: ils coïncident à
environ $\num{1e-5}$ en relatif --- l'approximation est bien meilleure
que toute mesure ici.

\textbf{6.} Par la \cref{prop:g12:light-as-wave:hair}~: un obstacle de
largeur $d$ diffracte comme une fente de largeur $d$, donc
$L = 2\lambda D/d$ tient encore.

\textbf{7.} $d = 2\lambda D/L = 2 \times \num{6.33e-7} \times
3.00/0.054 \approx \qty{70}{\micro m}$.

\textbf{8.} $L = \qty{5.2}{cm} \to d = \qty{73}{\micro m}$~;
$L = \qty{5.6}{cm} \to d = \qty{68}{\micro m}$~: donc
$d = \qty{70 \pm 3}{\micro m}$.

\textbf{9.} $d = \num{3.80e-6}/0.076 = \qty{50}{\micro m}$~: celui de
l'ami est plus fin --- et $L \propto 1/d$, donc le cheveu plus fin
jette le motif plus large.

\textbf{10.} Non~: une lampe est blanche (chaque
\omterm{def:g12:mechanical-waves:wavelength}{longueur d'onde} peint un motif différent, et ils se
brouillent) et incohérente-et-étendue (pas de faisceau unique propre à
diffracter). Le \omterm{def:g11:color-light-sources:lamps}{laser} est \omterm{def:g10:light-spectra:white}{monochromatique} et
directionnel.

\textbf{11.} Seules des copies d'une même \omterm{def:g10:signals-and-waves:wave}{onde} sont
\omterm{def:g12:light-as-wave:coherent}{cohérentes}~; deux sources indépendantes dérivent hors
phase et les franges s'effacent (\cref{rem:g12:light-as-wave:lamps}).

\textbf{12.} Au centre $\delta = 0 = 0 \times \lambda$ pour chaque
$\lambda$~: \omterm{prop:g12:light-as-wave:conditions}{constructive} quelle que soit la
\omterm{def:g12:mechanical-waves:wavelength}{longueur d'onde}.

\textbf{13.} $i = 4.26/10 = \qty{4.3}{mm}$~; un empan de dix se lit
avec la même erreur de règle qu'une seule \omterm{def:g12:light-as-wave:fringes}{interfrange}, donc
l'erreur sur $i$ est divisée par dix.

\textbf{14.} $\lambda = i\,a'/D = \num{4.26e-3} \times
\num{2.50e-4}/2.00 = \num{5.33e-7} \approx \qty{533}{nm}$.

\textbf{15.} $\qty{533}{nm}$ se trouve dans $\qty{532 \pm 10}{nm}$~:
cohérent~; cette \omterm{def:g12:mechanical-waves:wavelength}{longueur d'onde} est verte.

\textbf{16.} $\tan\theta = 0.43/1.00$, donc $\theta = 23^\circ$. Ici
$\tan\theta = 0.43$ tandis que $\sin\theta = 0.40$~: à $23^\circ$
l'identification à petit angle de $\sin$, $\tan$ et $\theta$ se trompe
de presque $10\%$ --- utiliser les fonctions exactes.

\textbf{17.} $d = \lambda/\sin\theta = \num{6.33e-7}/0.396 \approx
\qty{1.6}{\micro m}$.

\textbf{18.} $\tan\theta = 1.65$, $\theta = 59^\circ$,
$\sin\theta = 0.855$~: $d = \num{6.33e-7}/0.855 \approx
\qty{0.74}{\micro m}$.

\textbf{19.} $1.6/0.74 \approx 2.2$~: les pistes du DVD sont deux fois
plus denses, et ses creux sont correspondamment plus petits (lus avec un
\omterm{def:g11:color-light-sources:lamps}{laser} de \omterm{def:g12:mechanical-waves:wavelength}{longueur d'onde} plus courte) --- plusieurs
fois plus de données sur le même disque de \qty{12}{cm}.

\textbf{20.} Cheveu $\qty{70 \pm 3}{\micro m}$~; cheveu de l'ami
\qty{50}{\micro m}~; \omterm{def:g11:color-light-sources:lamps}{laser} vert \qty{533}{nm}~; pistes CD
\qty{1.6}{\micro m}~; pistes DVD \qty{0.74}{\micro m} --- une
\omterm{def:g12:mechanical-waves:wavelength}{longueur d'onde} connue transforme la géométrie de chaque
motif en une longueur, si bien que de la lumière d'un demi-micromètre
est une règle graduée au demi-micromètre.
\end{solution}
